% GENERATED FILE. Edit canonical lessons, then run npm run generate:books.
% canonical-source-hash: fnv1a64:2dd3374e22258f57
% canonical-lessons: FR-C377-reviser, FR-C377-revision, FR-C377-memoriser, FR-C377-par-coeur, FR-C377-manuel

\chapter{Avant l'examen}
\label{ch:fr-avant-l-examen}
\hlchaptermodality{377}

\begin{chapteropening}
\textbf{By the end of this chapter:} \emph{I can talk about getting ready for an exam: revising, revision, memorising, learning by heart and the textbook.}

Complete the last lesson of chapter 377: I can talk about getting ready for an exam: revising, revision, memorising, learning by heart and the textbook.
\end{chapteropening}

\section[réviser]{réviser — to revise, to study for (an exam); to service (a car)}

\label{lesson:FR-C377-reviser}



\begin{quote}
Before the new one: say the French for the semester, then the French for the lecture hall.
\end{quote}

\subsection*{You'll want to know: réviser}
\textbf{réviser} — \textquotedblleft{}to revise, to study for (an exam); to service (a car)\textquotedblright{}.

\begin{quote}
Je révise mes cours tous les soirs.
\end{quote}

\begin{quote}
\emph{I revise my notes every evening.}
\end{quote}

Belgium and Québec: \emph{étudier}. \emph{Faire réviser sa voiture} — to have the car serviced.

\begin{cousinweb}
\emph{réviser} is a learned borrowing from Latin \emph{revīsere} \textquotedblleft{}to view again, to revisit\textquotedblright{} (EtymDB; Lewis \& Short), built on \emph{vīsere}, from \emph{vidēre} \textquotedblleft{}to see\textquotedblright{}.
\end{cousinweb}

\begin{grammarlens}[title={What you've built}]
One more word for telling, explaining and describing at B1.
\end{grammarlens}

\subsection*{Your turn}
\begin{itemize}
\raggedright
  \item \emph{Say it:} \emph{réviser}
  \item \emph{Say it:} \emph{réviser}, once more
  \item \emph{Say it:} a sentence of your own with \emph{réviser}
  \item \emph{Recall:} say the French for the semester, then the French for the lecture hall, then say \emph{réviser} again
\end{itemize}

\begin{tcolorbox}[breakable,colback=teal!4,colframe=teal!35!black,title={Before you move on}]
What does réviser mean? (\textquotedblleft{}To revise, to study for (an exam); to service (a car)\textquotedblright{}.) Say it once more.
\end{tcolorbox}
\section[la révision]{la révision — the revision; the service (of a car)}

\label{lesson:FR-C377-revision}



\begin{quote}
Before the new one: say the French for the lecture hall, then the French for to revise.
\end{quote}

\subsection*{You'll want to know: la révision}
\textbf{la révision} — \textquotedblleft{}the revision; the service (of a car)\textquotedblright{}.

\begin{quote}
Pendant les révisions, je ne sors pas.
\end{quote}

\begin{quote}
\emph{During revision, I do not go out.}
\end{quote}

In the plural for exam revision: \emph{les révisions}.

\begin{cousinweb}
\emph{révision} is a learned borrowing from Latin \emph{revīsiō} \textquotedblleft{}a seeing again\textquotedblright{} (EtymDB), from \emph{revidēre} \textquotedblleft{}to go to see again\textquotedblright{}, built on \emph{vidēre} \textquotedblleft{}to see\textquotedblright{} (Lewis \& Short).
\end{cousinweb}

\begin{grammarlens}[title={What you've built}]
One more word for telling, explaining and describing at B1.
\end{grammarlens}

\subsection*{Your turn}
\begin{itemize}
\raggedright
  \item \emph{Say it:} \emph{la révision}
  \item \emph{Say it:} \emph{la révision}, once more
  \item \emph{Say it:} a sentence of your own with \emph{la révision}
  \item \emph{Recall:} say the French for the lecture hall, then the French for to revise, then say \emph{la révision} again
\end{itemize}

\begin{tcolorbox}[breakable,colback=teal!4,colframe=teal!35!black,title={Before you move on}]
What does la révision mean? (\textquotedblleft{}The revision; the service (of a car)\textquotedblright{}.) Say it once more.
\end{tcolorbox}
\section[mémoriser]{mémoriser — to memorise}

\label{lesson:FR-C377-memoriser}



\begin{quote}
Before the new one: say the French for to revise, then the French for the revision.
\end{quote}

\subsection*{You'll want to know: mémoriser}
\textbf{mémoriser} — \textquotedblleft{}to memorise\textquotedblright{}.

\begin{quote}
J'ai mémorisé tous les verbes.
\end{quote}

\begin{quote}
\emph{I memorised all the verbs.}
\end{quote}

From \emph{mémoire}, memory.

\begin{cousinweb}
\emph{mémoriser} is built on \emph{mémoire} \textquotedblleft{}memory\textquotedblright{}, from Latin \emph{memoria} \textquotedblleft{}memory\textquotedblright{}, from \emph{memor} \textquotedblleft{}mindful\textquotedblright{} (Lewis \& Short).
\end{cousinweb}

\begin{grammarlens}[title={What you've built}]
One more word for telling, explaining and describing at B1.
\end{grammarlens}

\subsection*{Your turn}
\begin{itemize}
\raggedright
  \item \emph{Say it:} \emph{mémoriser}
  \item \emph{Say it:} \emph{mémoriser}, once more
  \item \emph{Say it:} a sentence of your own with \emph{mémoriser}
  \item \emph{Recall:} say the French for to revise, then the French for the revision, then say \emph{mémoriser} again
\end{itemize}

\begin{tcolorbox}[breakable,colback=teal!4,colframe=teal!35!black,title={Before you move on}]
What does mémoriser mean? (\textquotedblleft{}To memorise\textquotedblright{}.) Say it once more.
\end{tcolorbox}
\section[par cœur]{par cœur — by heart}

\label{lesson:FR-C377-par-coeur}



\begin{quote}
Before the new one: say the French for the revision, then the French for to memorise.
\end{quote}

\subsection*{You'll want to know: par cœur}
\textbf{par cœur} — \textquotedblleft{}by heart\textquotedblright{}.

\begin{quote}
Il apprend le poème par cœur.
\end{quote}

\begin{quote}
\emph{He is learning the poem by heart.}
\end{quote}

\emph{Connaître par cœur} — to know by heart.

\begin{cousinweb}
\emph{cœur} is inherited from Latin \emph{cor} \textquotedblleft{}heart\textquotedblright{} (EtymDB; Lewis \& Short), once thought to be the seat of memory.
\end{cousinweb}

\begin{grammarlens}[title={What you've built}]
One more word for telling, explaining and describing at B1.
\end{grammarlens}

\subsection*{Your turn}
\begin{itemize}
\raggedright
  \item \emph{Say it:} \emph{par cœur}
  \item \emph{Say it:} \emph{par cœur}, once more
  \item \emph{Say it:} a sentence of your own with \emph{par cœur}
  \item \emph{Recall:} say the French for the revision, then the French for to memorise, then say \emph{par cœur} again
\end{itemize}

\begin{tcolorbox}[breakable,colback=teal!4,colframe=teal!35!black,title={Before you move on}]
What does par cœur mean? (\textquotedblleft{}By heart\textquotedblright{}.) Say it once more.
\end{tcolorbox}
\section[le manuel]{le manuel — the textbook, the manual}

\label{lesson:FR-C377-manuel}



\begin{quote}
Before the new one: say the French for to memorise, then the French for by heart.
\end{quote}

\subsection*{You'll want to know: le manuel}
\textbf{le manuel} — \textquotedblleft{}the textbook, the manual\textquotedblright{}.

\begin{quote}
Le manuel de maths coûte cher.
\end{quote}

\begin{quote}
\emph{The maths textbook is expensive.}
\end{quote}

As an adjective: done by hand, \emph{un travail manuel}.

\begin{cousinweb}
\emph{manuel} comes from Late Latin \emph{manuāle} \textquotedblleft{}a handbook\textquotedblright{} (EtymDB), from \emph{manuālis} \textquotedblleft{}of the hand\textquotedblright{}, built on \emph{manus} \textquotedblleft{}hand\textquotedblright{} (Lewis \& Short).
\end{cousinweb}

\begin{grammarlens}[title={What you've built}]
One more word for telling, explaining and describing at B1.
\end{grammarlens}

\subsection*{Your turn}
\begin{itemize}
\raggedright
  \item \emph{Say it:} \emph{le manuel}
  \item \emph{Say it:} \emph{le manuel}, once more
  \item \emph{Say it:} a sentence of your own with \emph{le manuel}
  \item \emph{Recall:} say the French for to memorise, then the French for by heart, then say \emph{le manuel} again
\end{itemize}

\begin{tcolorbox}[breakable,colback=teal!4,colframe=teal!35!black,title={Before you move on}]
What does le manuel mean? (\textquotedblleft{}The textbook, the manual\textquotedblright{}.) Say it once more.
\end{tcolorbox}
